\lysh{14184}{4}

\begin{alist}
\item Η συνάρτηση ορίζεται για τους πραγματικούς αριθμούς $x$ για τους οποίους ισχύει:
\begin{gather*}
1 + (x + 2)x \ne 0\text{, δηλαδή} \\
x^2 + 2x + 1 \ne 0\text{, οπότε} \\
(x + 1)^2 \ne 0\text{και τελικά} \\
x \ne -1.
\end{gather*}
Άρα $A_f = \mathbb{R} - \{-1\}$.

\item Έχουμε: $f(x) = \dfrac{(x + 2)(x + 1)^2}{1 + (x + 2)x} = \dfrac{(x + 2)(x + 1)^2}{(x + 1)^2} = x + 2$ και η γραφική της παράσταση
είναι ευθεία με εξίσωση $y = x + 2$ και $x \ne -1$.

\item 
\begin{rlist}
\item Η γραφική παράσταση της $f$ τέμνει τη γραφική παράσταση της $g$ σε σημεία των
οποίων οι τετμημένες είναι λύσεις της εξίσωσης: $x + 2 = x^2$, για $x \ne -1$.
Έχουμε λοιπόν $x + 2 = x^2 \Leftrightarrow x^2 - x - 2 = 0$ που είναι εξίσωση $2$ου βαθμού με ρίζες που
έχουν άθροισμα $-\dfrac{-1}{1} = 1$ και γινόμενο $\dfrac{-2}{1} = -2$. Άρα $x_1 = 2$ και $x_2 = -1$ (απορρίπτεται).
Άρα οι γραφικές παραστάσεις έχουν ένα κοινό σημείο, το $(2, 4)$, αφού $f(2) = g(2) = 4$.

\item H γραφική παράσταση της $f$ βρίσκεται πάνω από τη γραφική παράσταση της $g$ για
τις τιμές του $x$ οι οποίες είναι λύσεις της ανίσωσης:
\[x + 2 > x^2 \Leftrightarrow x^2 - x - 2 < 0.\]
Η ανίσωση είναι $2$ου βαθμού με $\alpha = 1 > 0$ και ρίζες $x_1 = 2$ και $x_2 = -1$, οπότε αληθεύει
για $x \in (-1, 2)$.
\end{rlist}
\end{alist}

